\section*{Chapter \ref{ch:m3:perpetual-futures} --- Perpetual Futures}

\begin{solution}{exo:m3:perpetual-futures:1}
The rate is positive, so the long pays the shorts $2 \times 60\,000 \times 0.03\% = \$36$.
\end{solution}

\begin{solution}{exo:m3:perpetual-futures:2}
0.02\% is inside the dead zone ($-0.04\%$ to $+0.06\%$): funding 0.01\%. 0.10\%: $0.10\% - 0.05\% = 0.05\%$.
$-0.08\%$: $-0.08\% + 0.05\% = -0.03\%$.
\end{solution}

\begin{solution}{exo:m3:perpetual-futures:3}
A single trade, possibly small or manipulated, could move the last price and trigger liquidations; the median of
an index-based funding estimate, an index-plus-average-basis estimate and the last trade ignores one outlier.
\end{solution}

\begin{solution}{exo:m3:perpetual-futures:4}
At 50\,000: $60\,000\,(1/60\,000 - 1/50\,000) = -0.2$ bitcoin, $-\$10\,000$. At 75\,000: $+0.2$ bitcoin,
$+\$15\,000$. In dollars the P\&L is linear, $60\,000\,(S/60\,000 - 1)$; in bitcoin it is not.
\end{solution}

\begin{solution}{exo:m3:perpetual-futures:5}
$10\,000 \times 0.000001 \times 300 = 3$ bitcoin, \$180\,000 at 60\,000; the dollar value would differ at another
bitcoin price although the ether move is the same.
\end{solution}

\begin{solution}{exo:m3:perpetual-futures:6}
$0.01\% \times 3 \times 365 = 10.95\%$: about USD~1.095 million.
\end{solution}

\begin{solution}{exo:m3:perpetual-futures:7}
12.93\% (2020), 23.20\% (2021), 2.95\% (2022), 5.79\% (2023), 8.89\% (2024), 3.70\% (2025) and 2.02\% (2026, at the
mean rate to September). Below 5\% in 2022, 2025 and 2026.
\end{solution}

\begin{solution}{exo:m3:perpetual-futures:8}
Price moves cancel but funding can turn negative, the short leg's margin can run out in a rally and be liquidated,
the venues holding both legs can fail, the \omterm{def:m3:blockchains-for-traders:stable}{stablecoin} margin can \omterm{def:m3:blockchains-for-traders:stable}{depeg}, and the basis at exit can differ from entry.
\end{solution}

\begin{solution}{pb:m3:perpetual-futures:1}
\textbf{1.} USD~11.92 million. \textbf{2.} $2 \times (0.10\% + 0.05\%) \times 100$ million: USD~0.30 million.
\textbf{3.} 10.95\% a year. \textbf{4.} The average premium spent most intervals inside the dead zone, where funding
equals the interest component. \textbf{5.} USD~4.16 million. \textbf{6.} $30\% + 0.5\% \times 1.3 = 30.65\%$ of
notional: USD~30.65 million. \textbf{7.} USD~130.65 million; $11.62/130.65 = 8.89\%$. \textbf{8.} 11.62\%.
\textbf{9.} The short is liquidated when its loss leaves only maintenance margin, at a rise of $(30.65 -
0.5)/100.5 = 30.0\%$; the fund is left long spot and unhedged, having sold nothing high. \textbf{10.} The spot sits
in a separate account; its gain is unrealised and not available as margin unless the venue counts it.
\textbf{11.} The short pays funding while hedged; the income becomes a cost. \textbf{12.} Both legs and the margin
sit on venues: their failure is a loss of up to the whole capital. \textbf{13.} The short's margin and P\&L are in
USDT; a \omterm{def:m3:blockchains-for-traders:stable}{depeg} changes their dollar value. \textbf{14.} If the perpetual trades cheap to spot at exit, closing the
short costs more than the spot sale returns relative to entry; the premium at entry and exit adds to the result.
\textbf{15.} The trade's income is highest just before crashes, when margins spike, liquidations cascade and funding
can reverse: its returns and its risks arrive together. \textbf{16.} Not strictly: it earns an uncertain stream and
bears margin, venue and \omterm{def:m3:blockchains-for-traders:stable}{stablecoin} risk; it is a carry trade with a hedge. \textbf{17.} Longs pay for leveraged
exposure without borrowing or holding spot; funding is the price of that leverage. \textbf{18.} By venue credit
limits, by each venue's depth and funding, and so that a rally's margin calls can be met on every venue at once.
\textbf{19.} \emph{Named result:} 8.89\% on capital over 2024, with 30.65\% of the notional as margin to survive a
30\% rally. \textbf{20.} The right to hold a leveraged long position without an expiry, for another eight hours.
\end{solution}

\begin{solution}{iq:m3:perpetual-futures:1}
A future with no expiry; holders exchange funding payments every interval, longs paying shorts when the contract is
rich to the index and the reverse when it is cheap, which rewards arbitrageurs for pushing it back.
\iqlookfor{funding as the convergence mechanism.}
\end{solution}

\begin{solution}{iq:m3:perpetual-futures:2}
The last trade is easy to move for a moment; liquidating on it would reward manipulation. The mark is built from the
index (capped, multi-venue) and a smoothed basis, for example the median of an index-plus-funding-basis price, an
index-plus-moving-average-basis price and the last trade.
\iqlookfor{manipulation resistance and the median.}
\end{solution}

\begin{solution}{iq:m3:perpetual-futures:3}
Linear: P\&L in the quote currency, the default for dollar-based traders. Inverse: dollar notional, coin margin and
settlement; P\&L in coin is concave, and a short plus coin margin is a dollar-stable position, useful to coin-based
holders such as miners. Quanto: exposure to one asset settled in another at a fixed rate, adding the settlement asset
and the correlation to the risk.
\iqlookfor{currency of P\&L and who naturally holds each.}
\end{solution}

\begin{solution}{iq:m3:perpetual-futures:4}
Because the average premium sits inside the dead zone, where the clamp pins funding at the interest component of
0.01\% whatever the premium, as long as it stays within 0.05\% of it.
\iqlookfor{the clamp.}
\end{solution}

\begin{solution}{iq:m3:perpetual-futures:5}
Negative funding; margin calls on the short legs in a rally, possibly on several venues at once, and liquidation if
unmet; venue failure; \omterm{def:m3:blockchains-for-traders:stable}{stablecoin} \omterm{def:m3:blockchains-for-traders:stable}{depeg}; basis moves at unwind; operational risk in moving collateral. Limits per
venue, margin buffers, cross-margin where possible, and a plan to unwind.
\iqlookfor{liquidity of margin as the main risk.}
\end{solution}

\begin{solution}{iq:m3:perpetual-futures:6}
Two hours before, three quarters of the samples are in and, with linear weights, they already carry about 56\% of the
weight. Reproduce the venue's premium from its book (impact notional, 5-second sampling), compute the weighted
average so far, forecast the remaining samples (their current level, mean reversion) and apply the clamp and caps;
the dead zone makes the answer exact whenever the forecast stays inside it.
\iqlookfor{reproduce the index and the weights, then forecast what remains.}
\end{solution}
